% ch04-histories.tex
\chapter{Histories and Continuation Languages}\label{ch:histories}
\spec{Which properties belong to a current state, and which belong to the history that produced it?}
{Finite words and the prefix order; deterministic transition systems; the Myhill--Nerode construction; safety and liveness (Alpern and Schneider).}
{The set of admissible extensions of the account determines its balance on viable histories (Result~\ref{prop:admissible-balance}).}
{None.}
{Chapters~\ref{ch:common} and~\ref{ch:quotients}.}

Up to this point the cases have been arbitrary. From now on they are histories: finite sequences of events, ordered by the relation of one being the beginning of another. This chapter sets out the elementary structure of such sets of histories, introduces the transition systems that generate them, and draws the distinction on which the rest of the volume depends. A property that can be read from the state a system occupies is a property of a particular representation of its history, and whether a given property can be read from a given state is exactly the factorization question of Chapter~\ref{ch:quotients}. The chapter closes with the classical separation of properties into those whose violation is detected on a finite history and those that remain open at every finite history, which is the form in which obligations will appear later.

\section{Words, prefixes, and history sets}

\begin{definition}[Words and the prefix order]
Let $\Events$ be a nonempty set of events. $\Events^{*}$ denotes the set of finite words over $\Events$ with concatenation, empty word $\varepsilon$, and length $|w|$. Write $u\prefixeq v$ if $v=uw$ for some word $w$.
\end{definition}

\begin{established}[Structure of the prefix order]\label{prop:prefix}
The relation $\prefixeq$ is a partial order on $\Events^{*}$. The prefixes of a word $v$ form a chain with exactly $|v|+1$ elements, one of each length from $0$ to $|v|$, and every nonempty word has a unique immediate predecessor, obtained by deleting its last event.
\end{established}
\source{Standard.}
\begin{proof}
Reflexivity holds with $w=\varepsilon$ and transitivity by associativity of concatenation. If $u\prefixeq v$ and $v\prefixeq u$ then $|u|=|v|$ and $v=uw$ forces $w=\varepsilon$, so $u=v$. A prefix of $v$ of length $\ell$ is determined by $\ell$ and consists of the first $\ell$ events of $v$, which gives the chain. The last statement is the case $\ell=|v|-1$.
\end{proof}

\begin{definition}[History set and extensions]\label{def:history-set}
A \emph{history set} is a nonempty set $\Hist\subseteq\Events^{*}$ that is closed under prefixes. For $h\in\Hist$, the set of \emph{extensions} of $h$ is $\Ext(h)=\{w\in\Events^{*}: hw\in\Hist\}$.
\end{definition}

\begin{established}[Extensions]\label{prop:extensions}
For $h\in\Hist$ the set $\Ext(h)$ is itself a history set. If $u\in\Ext(h)$, then $\Ext(hu)=\{w : uw\in\Ext(h)\}$.
\end{established}
\source{Standard.}
\begin{proof}
We have $\varepsilon\in\Ext(h)$ since $h\in\Hist$. If $w_1w_2\in\Ext(h)$ then $hw_1w_2\in\Hist$, so $hw_1\in\Hist$ by closure under prefixes, and $w_1\in\Ext(h)$. For the second statement, $w\in\Ext(hu)$ means $huw\in\Hist$, which means $uw\in\Ext(h)$.
\end{proof}

A history set is the same thing as a rooted tree whose edges are labelled by events, with no two edges leaving a node carrying the same label: nodes are histories, the root is $\varepsilon$, and $h$ is joined to $ha$ whenever both lie in $\Hist$. The set $\Ext(h)$ is the subtree rooted at $h$. The constraints of a system, whether physical or normative, show up in this picture as missing branches.

\section{Transition systems and the histories they generate}

\begin{definition}[Deterministic transition system]\label{def:lts}
A \emph{deterministic labelled transition system} is a tuple $T=(S,s_0,\Events,\delta)$ with a set $S$ of states, an initial state $s_0\in S$, and a partial function $\delta:S\times\Events\rightharpoonup S$. Extend $\delta$ to words by $\delta^{*}(s,\varepsilon)=s$ and $\delta^{*}(s,wa)=\delta(\delta^{*}(s,w),a)$ whenever both sides are defined, and leave $\delta^{*}$ undefined otherwise. The set of histories generated by $T$ is $\Hist(T)=\{w: \delta^{*}(s_0,w)\text{ is defined}\}$, and the \emph{state map} is $s_T:\Hist(T)\to S$, $h\mapsto\delta^{*}(s_0,h)$.
\end{definition}

\begin{established}[Histories and transition systems]\label{prop:lts}
\begin{enumerate}[nosep,label=(\alph*)]
  \item $\Hist(T)$ is a history set for every $T$.
  \item Every history set $\Hist$ is generated by the \emph{tree system} $T_{\Hist}=(\Hist,\varepsilon,\Events,\delta)$ with $\delta(h,a)=ha$ when $ha\in\Hist$.
  \item $\Hist$ is a regular language if and only if $\Hist=\Hist(T)$ for some $T$ with finitely many states.
\end{enumerate}
\end{established}
\source{Standard automata theory; see Hopcroft and Ullman, \emph{Introduction to Automata Theory, Languages, and Computation}.}
\begin{proof}
(a) If $\delta^{*}(s_0,wa)$ is defined then $\delta^{*}(s_0,w)$ is defined, so $\Hist(T)$ is prefix-closed, and it contains $\varepsilon$. (b) By induction on length, $\delta^{*}(\varepsilon,w)=w$ exactly when $w\in\Hist$. (c) A finite $T$ becomes a complete deterministic automaton on $S\cup\{\bot\}$ by sending every undefined transition to a new sink $\bot$ and declaring every state of $S$ accepting; this accepts $\Hist(T)$. Conversely let $\Hist$ be regular and let $M$ be its minimal complete deterministic automaton. A state of $M$ reached by a word of $\Hist$ is accepting. A state reached by a word outside $\Hist$ has empty residual language, because $\Hist$ is prefix-closed, so all such states coincide with a single sink. Deleting the sink leaves a finite partial system in which every state is reached by a word of $\Hist$ and which generates $\Hist$.
\end{proof}

The tree system is the largest and least informative way of presenting a history set, in that its states are the histories themselves. Other systems that generate the same history set carry less, and the whole question of this volume is which of them carries enough.

\section{Properties of states and properties of histories}

\begin{definition}[Property of histories, state property]\label{def:state-property}
A \emph{property} of a history set $\Hist$ is a function $P:\Hist\to D$ into a set of verdicts. For a transition system $T$ with $\Hist(T)=\Hist$, the property $P$ is a \emph{state property of $T$} if there is a function $p:s_T(\Hist)\to D$ with $P=p\circ s_T$.
\end{definition}

\begin{established}[State properties are factorizations]\label{prop:state-prop}
Let $T$ generate $\Hist$ and let $P$ be a property of $\Hist$. Then $P$ is a state property of $T$ if and only if $\ker s_T\subseteq\ker P$.
\end{established}
\source{Result~\ref{thm:factorization} with $r=s_T$ and $f=P$.}

Whether a property is a state property is therefore not a feature of the property alone. It is a relation between the property and a way of presenting the history set.

\begin{example}[Two presentations of the account]\label{ex:two-presentations}
Let $\Hist=\Events^{*}$ for the events of Example~\ref{ex:bank}. Let $T_1$ have state set $\ZZ$, initial state $0$, and transitions $\delta(n,\mathrm{d}(k))=n+k$ and $\delta(n,\mathrm{w}(k))=n-k$. Let $T_2$ have state set $\ZZ\times\{0,1\}$, initial state $(0,0)$, and transitions $\delta((n,f),\mathrm{d}(k))=(n+k,f)$ and $\delta((n,f),\mathrm{w}(k))=(n-k,\max(f,[n-k<0]))$, where $[\cdot]$ is $1$ when the condition holds and $0$ otherwise. Then $s_{T_1}=\beta$ and $s_{T_2}=(\beta,q_{\mathrm{neg}})$. For $T_1$ this is immediate. For $T_2$ the claim follows by induction on length: it holds for $\varepsilon$ since $q_{\mathrm{neg}}(\varepsilon)=0$; a deposit changes the balance of the history to $\beta(h)+k$, and if that is negative then $\beta(h)<0$, so $q_{\mathrm{neg}}(h)=1$ already and the flag is unchanged; a withdrawal changes the balance to $\beta(h)-k$, and the new prefix contributes $[\beta(h)-k<0]$ to the flag. By Result~\ref{prop:state-prop}, the balance is a state property of both systems and $q_{\mathrm{neg}}$ is a state property of $T_2$ only, since the balance does not answer $q_{\mathrm{neg}}$.
\end{example}

\begin{example}[Document]
For the document of Example~\ref{ex:document}, the system whose state is the current text $\tau(h)\in C^{m}$ makes ``what does section $j$ contain now?'' a state property and does not make ``has section $j$ ever contained $c$?'' one. A system whose state is the entire history, the tree system, makes every property of histories a state property and is therefore the extreme case.
\end{example}

The word \emph{state} in the common usage of a system is therefore relative to a presentation that someone chose, and a state description is a representation of histories in the sense of Chapter~\ref{ch:common}. The question of which presentation is sufficient is the sufficiency question that Part~II treats in full.

\subsection*{Continuation equivalence}

A state map has a further property that an arbitrary representation lacks: it is closed under the passage of time. Two histories that reach the same state remain at the same state after every common sequence of further events. The following construction isolates what this requires of a property.

\begin{definition}[Residual and continuation equivalence]\label{def:continuation-equiv}
For a property $P$ of $\Hist$ and a history $h$, the \emph{residual} is $P_h:\Ext(h)\to D$, $P_h(w)=P(hw)$. Histories $h$ and $h'$ are \emph{continuation equivalent for $P$}, written $h\approx_P h'$, if $\Ext(h)=\Ext(h')$ and $P_h=P_{h'}$.
\end{definition}

\begin{established}[The Myhill--Nerode construction]\label{prop:nerode}
Let $P$ be a property of $\Hist$.
\begin{enumerate}[nosep,label=(\alph*)]
  \item $\approx_P$ is an equivalence relation contained in $\ker P$.
  \item If $h\approx_P h'$ and $u\in\Ext(h)$, then $hu\approx_P h'u$.
  \item If $\sim$ is an equivalence relation on $\Hist$ contained in $\ker P$ such that $h\sim h'$ implies $\Ext(h)=\Ext(h')$ and $hu\sim h'u$ for all $u\in\Ext(h)$, then ${\sim}\subseteq{\approx_P}$.
\end{enumerate}
\end{established}
\source{Myhill (1957) and Nerode (1958), in the formulation for properties of histories.}
\begin{proof}
(a) Reflexivity, symmetry, and transitivity hold because equality of residuals does. Taking $w=\varepsilon$ in $P_h=P_{h'}$ gives $P(h)=P(h')$. (b) By Result~\ref{prop:extensions}, $\Ext(hu)=\{w: uw\in\Ext(h)\}=\Ext(h'u)$, and for such $w$ we have $P_{hu}(w)=P(huw)=P_h(uw)=P_{h'}(uw)=P_{h'u}(w)$. (c) Let $h\sim h'$ and $w\in\Ext(h)$. Then $hw\sim h'w$, so $P(hw)=P(h'w)$ because ${\sim}\subseteq\ker P$. Together with $\Ext(h)=\Ext(h')$ this gives $P_h=P_{h'}$.
\end{proof}

\begin{established}[Sufficiency of state, and minimality of the Nerode system]\label{prop:nerode-system}
Let $P$ be a property of $\Hist$.
\begin{enumerate}[nosep,label=(\alph*)]
  \item The system $T_P=(\Hist/{\approx_P},[\varepsilon],\Events,\delta)$ with $\delta([h],a)=[ha]$ for $ha\in\Hist$ is well defined, generates $\Hist$, and has $P$ as a state property.
  \item If $T$ is any transition system generating $\Hist$ for which $P$ is a state property, then $\ker s_T\subseteq{\approx_P}$. In particular $T$ has at least $|\Hist/{\approx_P}|$ reachable states.
\end{enumerate}
\end{established}
\source{Myhill and Nerode; see Hopcroft and Ullman.}
\begin{proof}
(a) The transition is well defined by Result~\ref{prop:nerode}(b) with $u=a$, and $\Ext$ is constant on classes by definition of $\approx_P$. By induction $s_{T_P}(h)=[h]$, so $T_P$ generates $\Hist$, and $P$ is a state property because ${\approx_P}\subseteq\ker P$. (b) Suppose $s_T(h)=s_T(h')=s$. Then $\Ext(h)=\{w:\delta^{*}(s,w)\text{ is defined}\}=\Ext(h')$, and for $w\in\Ext(h)$ we have $s_T(hw)=\delta^{*}(s,w)=s_T(h'w)$, so $P(hw)=P(h'w)$ because $P$ is a state property of $T$. Hence $h\approx_P h'$. The count follows from Result~\ref{prop:coarsening}.
\end{proof}

The residual construction expresses the difference between answering a question about the present history and answering it after every further event. A representation $r$ answers $P$ now if $\ker r\subseteq\ker P$. A representation that is the state map of a transition system and answers $P$ must satisfy the much stronger condition $\ker r\subseteq{\approx_P}$. For the account with $P=q_{\mathrm{neg}}$, the relation $\approx_P$ is the kernel $\ker q^{+}$ of Result~\ref{prop:account}, the Nerode system has the states $\NN\cup\{\bot\}$ with $[b]\xrightarrow{\mathrm{d}(k)}[b+k]$, $[b]\xrightarrow{\mathrm{w}(k)}[b-k]$ for $k\le b$, $[b]\xrightarrow{\mathrm{w}(k)}\bot$ for $k>b$, and $\bot$ absorbing. Chapter~\ref{ch:quotients} showed that the single flag $q_{\mathrm{neg}}$ answers the question about the present history while the Nerode system needs the balance for as long as no overdraft has occurred. Chapter~\ref{ch:incremental} takes up this contrast systematically.

\section{Safety and liveness}

Properties of histories arise in practice from properties of unbounded behaviour. A rule such as ``the balance never goes negative'' or ``every request is eventually answered'' speaks about all prefixes, or about the whole future, of a run. This section states, with proofs, the classical vocabulary for sorting such rules by what a finite history can reveal about them.

Let $\Events^{\omega}$ be the set of infinite words over $\Events$, and for the remainder of the section take the history set to be all of $\Events^{*}$. A \emph{behavioural property} is a set $\Pi\subseteq\Events^{\omega}$. For a finite word $u$ write $\sat_\Pi(u)$ if some $\sigma\in\Events^{\omega}$ has $u\sigma\in\Pi$, so that $u$ is still compatible with satisfying $\Pi$.

\begin{definition}[Safety and liveness]\label{def:safety-liveness}
A property $\Pi$ is a \emph{safety property} if for every $\sigma\notin\Pi$ some finite prefix $u$ of $\sigma$ has $\neg\sat_\Pi(u)$. It is a \emph{liveness property} if $\sat_\Pi(u)$ holds for every finite word $u$.
\end{definition}

A safety property is one whose violation is detected on a finite history that no continuation can repair; a liveness property is one for which no finite history is ever decisive.

\begin{established}[Closure characterization of safety]\label{prop:safety-closure}
Let
\[
\Pi_{\mathrm{cl}}=\{\sigma\in\Events^{\omega}: \sat_\Pi(u)\text{ for every finite prefix }u\text{ of }\sigma\}.
\]
Then $\Pi\subseteq\Pi_{\mathrm{cl}}$ for every $\Pi$, and $\Pi$ is a safety property if and only if $\Pi=\Pi_{\mathrm{cl}}$.
\end{established}
\source{Alpern and Schneider, ``Defining liveness,'' \emph{Information Processing Letters} 21 (1985).}
\begin{proof}
If $\sigma\in\Pi$ then every prefix $u$ of $\sigma$ satisfies $\sat_\Pi(u)$, as shown by $\sigma$ itself, so $\Pi\subseteq\Pi_{\mathrm{cl}}$. If $\Pi=\Pi_{\mathrm{cl}}$ and $\sigma\notin\Pi$, some prefix $u$ fails $\sat_\Pi$, so $\Pi$ is safe. If $\Pi$ is safe and $\sigma\in\Pi_{\mathrm{cl}}\setminus\Pi$, the definition of safety gives a prefix $u$ with $\neg\sat_\Pi(u)$, contradicting $\sigma\in\Pi_{\mathrm{cl}}$.
\end{proof}

\begin{established}[Decomposition]\label{thm:alpern-schneider}
Every property $\Pi$ is the intersection of a safety property and a liveness property.
\end{established}
\source{Alpern and Schneider (1985).}
\begin{proof}
Take $\Pi_S=\Pi_{\mathrm{cl}}$ and $\Pi_L=\Pi\cup(\Events^{\omega}\setminus\Pi_{\mathrm{cl}})$. The set $\Pi_{\mathrm{cl}}$ is safe: if $\sigma\notin\Pi_{\mathrm{cl}}$ then some prefix $u$ has $\neg\sat_\Pi(u)$, and no infinite extension of $u$ lies in $\Pi_{\mathrm{cl}}$, since each such extension has $u$ as a prefix. The set $\Pi_L$ is live: if $\sat_\Pi(u)$ then some extension of $u$ lies in $\Pi\subseteq\Pi_L$, and if $\neg\sat_\Pi(u)$ then every extension of $u$ lies outside $\Pi_{\mathrm{cl}}$ and so inside $\Pi_L$. Finally $\Pi_S\cap\Pi_L=(\Pi_{\mathrm{cl}}\cap\Pi)\cup\emptyset=\Pi$, using $\Pi\subseteq\Pi_{\mathrm{cl}}$.
\end{proof}

The set of finite histories compatible with a property is a history set in its own right.

\begin{established}[Viable histories]\label{prop:viable}
For nonempty $\Pi$ let $\Viable(\Pi)=\{u\in\Events^{*}:\sat_\Pi(u)\}$. Then $\Viable(\Pi)$ is a history set which is \emph{pruned}: every $u\in\Viable(\Pi)$ has some $a\in\Events$ with $ua\in\Viable(\Pi)$. A word outside $\Viable(\Pi)$ has all its extensions outside $\Viable(\Pi)$.
\end{established}
\source{Standard.}
\begin{proof}
Nonemptiness of $\Pi$ gives $\varepsilon\in\Viable(\Pi)$. If $uv$ is viable, with extension $\sigma$, then $u$ is viable, with extension $v\sigma$, which shows prefix-closure. If $u$ is viable via $\sigma=a\sigma'$ then $ua$ is viable via $\sigma'$, which shows pruning. The final claim is the contrapositive of prefix-closure.
\end{proof}

\begin{established}[Safety properties are pruned history sets]\label{prop:safety-trees}
For a pruned history set $V$ let $\Pi_V=\{\sigma:\text{every finite prefix of }\sigma\text{ lies in }V\}$. The maps $\Pi\mapsto\Viable(\Pi)$ and $V\mapsto\Pi_V$ are mutually inverse bijections between nonempty safety properties and pruned history sets.
\end{established}
\source{The correspondence between closed subsets of $\Events^{\omega}$ and pruned trees; see Kechris, \emph{Classical Descriptive Set Theory}.}
\begin{proof}
Let $V$ be pruned and $u\in V$. Starting from $u_0=u$, pruning gives events $a_i$ with $u_{i+1}=u_ia_i\in V$ for all $i$ (by dependent choice), and the limit $\sigma$ has all of its prefixes in $V$ because every prefix is a prefix of $u$ or equals some $u_i$. Thus $\sigma\in\Pi_V$ extends $u$ and $u\in\Viable(\Pi_V)$. If $u\notin V$ then no extension of $u$ lies in $\Pi_V$, so $u\notin\Viable(\Pi_V)$. Hence $\Viable(\Pi_V)=V$. This also shows $\Pi_V$ is safe, because $(\Pi_V)_{\mathrm{cl}}$ consists of the words all of whose prefixes lie in $\Viable(\Pi_V)=V$, which is $\Pi_V$ itself, and Result~\ref{prop:safety-closure} applies. For a nonempty safety property $\Pi$, Result~\ref{prop:viable} shows $\Viable(\Pi)$ is a pruned history set, and $\Pi_{\Viable(\Pi)}=\Pi_{\mathrm{cl}}=\Pi$ by Result~\ref{prop:safety-closure}.
\end{proof}

Thus a safety property is nothing other than a pruned history set read as a permission to continue: the property holds of a run exactly when the run never leaves the set. The histories in the set are the ones from which the property can still be met, and this is the finite-history form of the viability condition studied in the second volume. A liveness property has $\Viable(\Pi)=\Events^{*}$, since by definition every finite word is compatible with it. Violating it is a feature of the infinite run and never of a finite history.

\section{Admissible extensions and the memory they require}

\begin{definition}[Admissible extensions]\label{def:admissible-ext}
For a nonempty property $\Pi$ and a viable history $h$, the set of \emph{admissible extensions} of $h$ is $\Ext_\Pi(h)=\{w\in\Events^{*}: hw\in\Viable(\Pi)\}$.
\end{definition}

\begin{example}[The account that must not be overdrawn]\label{ex:no-overdraft}
In the account of Example~\ref{ex:bank} let $\Pi_{\mathrm{neg}}$ be the set of infinite runs all of whose finite prefixes have nonnegative balance. The set $V=\{h: q_{\mathrm{neg}}(h)=0\}$ is a history set, since the prefixes of a history with no negative prefix have none, and it is pruned, since a deposit of $1$ cannot create a negative prefix. Then $\Pi_{\mathrm{neg}}=\Pi_V$ and $\Viable(\Pi_{\mathrm{neg}})=V$ by Result~\ref{prop:safety-trees}. For $h\in V$ the admissible extensions are $\Ext_{\Pi}(h)=\{w:\beta(h)+m(w)\ge0\}$, where $m(w)\le0$ is the least balance among the prefixes of $w$.
\end{example}

\begin{newresult}[Admissible extensions determine the balance]\label{prop:admissible-balance}
For $h,h'\in V$ in Example~\ref{ex:no-overdraft}, $\Ext_{\Pi}(h)=\Ext_{\Pi}(h')$ if and only if $\beta(h)=\beta(h')$.
\end{newresult}
\assumes{Events, balances, and the property $\Pi_{\mathrm{neg}}$ are as in Example~\ref{ex:no-overdraft}; amounts are positive integers.}
\begin{proof}
If $\beta(h)=\beta(h')$ the descriptions of the two sets in Example~\ref{ex:no-overdraft} coincide. Suppose $\beta(h)=b<b'=\beta(h')$. The one-event word $w=\mathrm{w}(b')$ has $m(w)=-b'$ and $b'\ge1$. Then $b'+m(w)=0\ge0$, so $w\in\Ext_\Pi(h')$, while $b+m(w)=b-b'<0$, so $w\notin\Ext_\Pi(h)$.
\end{proof}

The result says that on viable histories the information needed to decide which transitions are admissible is the balance and nothing coarser, and combined with Result~\ref{prop:account} it identifies the least memory that keeps admissibility computable: the balance on the viable region and a single absorbing mark beyond it. This is the first instance of a pattern that organizes Part~IV: the representation required to adjudicate admissibility is the continuation quotient of the property that defines it.


\section{What finite histories can see of a behavioural property}\label{sec:closure-view}

The decomposition into safety and liveness has a consequence for retention that is easy to miss. A decision procedure that consults only finite histories cannot tell a behavioural property from its safety closure, and the part of a requirement that is lost in this way is exactly the liveness component.

\begin{newresult}[Finite histories see only the closure]\label{prop:closure-view}
For every nonempty $\Pi\subseteq\Events^{\omega}$, $\Viable(\Pi)=\Viable(\Pi_{\mathrm{cl}})$ and $\Pi_{\mathrm{cl}}=\Pi_{\Viable(\Pi)}$. Consequently every function of finite histories that is determined by viability, in particular the sets $\Ext_\Pi(h)$ of admissible extensions, is the same for $\Pi$ and for its closure.
\end{newresult}
\assumes{$\Pi$ nonempty; the history set is all of $\Events^{*}$.}
\begin{proof}
Since $\Pi\subseteq\Pi_{\mathrm{cl}}$, every word compatible with $\Pi$ is compatible with $\Pi_{\mathrm{cl}}$. Conversely, if $\sigma\in\Pi_{\mathrm{cl}}$ extends $u$, then $u$ is a prefix of $\sigma$ and $\sat_\Pi(u)$ holds by the definition of $\Pi_{\mathrm{cl}}$. So $\Viable(\Pi)=\Viable(\Pi_{\mathrm{cl}})$. The second equation is the definition of $\Pi_{\mathrm{cl}}$ read through $\Viable$: an infinite word belongs to $\Pi_{\mathrm{cl}}$ exactly when all of its prefixes are viable.
\end{proof}

\begin{example}[A requirement with a liveness component]\label{ex:live-component}
For the monitor of Example~\ref{ex:monitor} let $\Pi$ be the set of runs in which no prefix has three alarms outstanding and in which $\mathrm{reset}$ occurs infinitely often. The closure $\Pi_{\mathrm{cl}}$ is the set of runs with no prefix having three alarms outstanding: a finite word with that property extends to a member of $\Pi$ by appending resets forever, so it is compatible with $\Pi$, and a word with a prefix of three outstanding alarms is compatible with nothing. Hence $\Viable(\Pi)$ is the set of finite histories with $q_{\mathrm{ever}}=0$, and the decomposition of Result~\ref{thm:alpern-schneider} is $\Pi=\Pi_{\mathrm{cl}}\cap\Pi_L$ with $\Pi_L=\Pi\cup(\Events^{\omega}\setminus\Pi_{\mathrm{cl}})$. The property $\Pi$ is neither safe nor live. The finite history $\mathrm{alarm}\,\mathrm{alarm}\,\mathrm{tick}^{n}$ is viable for every $n$, however long the system goes without a reset, and no finite history ever shows the requirement of infinitely many resets to have failed. A memory sized to adjudicate $\Pi$ on finite histories is the Nerode system of $q_{\mathrm{ever}}$, with four states; the infinitary requirement costs nothing in memory because it cannot be adjudicated by memory at all.
\end{example}

The practical reading is that obligations of the form ``eventually'' are not retained as state. They constrain what the system must still do, and Chapter~\ref{ch:obligations} converts them into bounded-response obligations whose pending deadlines are state, which is a different object.

\section{Why the extension clause cannot be dropped}\label{sec:ext-clause}

Definition~\ref{def:continuation-equiv} requires $\Ext(h)=\Ext(h')$ in addition to $P_h=P_{h'}$. The clause looks like a technicality, and a short example shows that it is not.

\begin{example}[Histories that agree where both are defined]\label{ex:ext-needed}
Let $\Events=\{a,b\}$ and $\Hist=\{\varepsilon,a,b,bb\}$, and let $P$ be constant. Suppose we declared $h$ and $h'$ equivalent whenever $P_h$ and $P_{h'}$ agree on $\Ext(h)\cap\Ext(h')$. Then all four histories would be equivalent, and the resulting one-state system would have to allow $a$ and $b$ at its only state, hence generate a larger history set than $\Hist$. With the clause the classes are $\{\varepsilon\}$, $\{b\}$ and $\{a,bb\}$, since $\Ext(\varepsilon)=\{\varepsilon,a,b,bb\}$, $\Ext(b)=\{\varepsilon,b\}$ and $\Ext(a)=\Ext(bb)=\{\varepsilon\}$. The Nerode system has three states, an initial state with transitions $a$ to the end state and $b$ to a middle state, and a middle state with a transition $b$ to the end state, and it generates $\Hist$ exactly. No system generating $\Hist$ has fewer than three reachable states, by Result~\ref{prop:nerode-system}(b).
\end{example}

The example makes a point about what a property can be said to distinguish. A constant property distinguishes nothing about verdicts, yet the histories $a$, $b$ and $\varepsilon$ differ in how much future they have, and a system that is to generate the history set must retain that difference. In the common model this is the physical component: $\Reach(\phys(h))$ is part of what must be retained whatever questions are asked, and the obligations come on top of it.

\section*{Exercises}
\addcontentsline{toc}{section}{Exercises}

\begin{exercise}
Construct a transition system with four states that generates all words over $\{\mathrm{tick},\mathrm{alarm},\mathrm{reset}\}$ and for which ``at least three alarms are outstanding'' (the question of Example~\ref{ex:monitor}) is a state property.
\end{exercise}

\begin{exercise}
For your system, prove that $s_T=\rho$, and use Result~\ref{prop:state-prop} to conclude that $\ker\rho$ is contained in the kernel of the question. Then use Result~\ref{prop:nerode-system} to show that no transition system generating the same history set and having the question as a state property has fewer than four reachable states.
\end{exercise}

\begin{exercise}
Show that ``three alarms have ever been outstanding'' is not a state property of your system, by exhibiting two histories with the same state and different verdicts. Then compute the continuation equivalence for this second property and show that its Nerode system has exactly four states, namely the capped counts $0,1,2$ together with one absorbing state, and compare this with the four states of the first property.
\end{exercise}

\begin{exercise}
The threshold of the question is revised from three alarms to $j$ alarms. Compute $\approx_P$ for the revised question and show that the Nerode system has exactly $j+1$ states. Describe in words what is lost if the memory is kept at $\rho$ with cap $3$ and the threshold is revised to $5$.
\end{exercise}

\begin{exercise}
Classify the following behavioural properties of the monitor as safety, liveness, or neither: ``no prefix of the run has three alarms outstanding'' and ``some prefix of the run has three alarms outstanding''. For each, describe $\Viable(\Pi)$ and compute the decomposition of Result~\ref{thm:alpern-schneider}. Explain why the two properties are complements of one another and yet fall on opposite sides of the classification.
\end{exercise}

\begin{exercise}
Prove that for $\Pi$ a safety property and $h$ a viable history, $\Ext_\Pi(h)$ is a pruned history set, and that $\Ext_\Pi(hu)=\{w: uw\in\Ext_\Pi(h)\}$ for $u\in\Ext_\Pi(h)$. Use this to show that the map $h\mapsto\Ext_\Pi(h)$ on viable histories is constant on the classes of the continuation equivalence for the indicator of $\Viable(\Pi)$.
\end{exercise}


\begin{exercise}
For the property $\Pi$ of Example~\ref{ex:live-component}, compute $\Pi_{\mathrm{cl}}$ and $\Pi_L$ explicitly, verify that $\Pi_L$ is live, and exhibit a run in $\Pi_L\setminus\Pi$ that lies in $\Pi_{\mathrm{cl}}$ and one that lies outside $\Pi_{\mathrm{cl}}$.
\end{exercise}

\begin{exercise}
Prove that $\Pi$ is a liveness property if and only if $\Pi_{\mathrm{cl}}=\Events^{\omega}$. Deduce that a property that is both safe and live is all of $\Events^{\omega}$.\label{ex:live-closure}
\end{exercise}

\begin{exercise}
Change the history set of Example~\ref{ex:ext-needed} to $\{\varepsilon,a,b,bb,bbb\}$ with $P(h)=1$ exactly when $h$ contains an $a$. Compute the classes of $\approx_P$ and show that the Nerode system has five states, so that it coincides with the tree system. Explain which histories the extension clause separates that the verdicts alone would not.
\end{exercise}

